\chapter{Đi tiếp cùng em}
\markboth{Đi tiếp cùng em}{Keep Walking With You}


\section{Thầy không chỉ cứu một buổi học — thầy muốn cứu luôn lòng tin học tập của em.}
\begin{truyen}{Thầy không chỉ cứu một buổi học — thầy muốn cứu luôn lòn...}{Protect long-term hope.}
\chuhoa{C}{ó câu nói giúp học sinh qua được hôm nay; cũng có câu giữ các em còn tin vào việc học thêm nhiều tháng nữa. Người thầy tốt luôn nghĩ đến cả hai.}
\textit{(Some sentences help students survive today; others keep them believing in learning for many more months. A good teacher thinks about both.)}
\end{truyen}

\begin{baihoc}
Giữ được lòng tin học tập quý hơn kéo được một kết quả ngắn hạn.
\textit{(Preserving a student's belief in learning matters more than winning one short-term result.)}
\end{baihoc}

\begin{ghinhoanh}
\textbf{Short English to remember:}

Protect the learner, not just the lesson.

\end{ghinhoanh}

\begin{tuhocnhanh}
1. Điều gì làm em mất lòng tin vào việc học? \textit{(What makes you lose faith in learning?)}

2. Điều gì giúp em lấy lại niềm tin đó? \textit{(What helps you regain that faith?)}
\end{tuhocnhanh}
\ngancach


\section{Có những em không cần thêm áp lực — các em cần thêm người đi cùng.}
\begin{truyen}{Có những em không cần thêm áp lực — các em cần thêm ngườ...}{Companionship matters.}
\chuhoa{N}{hiều học sinh không yếu ở năng lực, mà yếu ở cảm giác đơn độc. Một người đi cùng đúng lúc có thể khiến các em không bỏ cuộc giữa đường.}
\textit{(Many students are not weak in ability, but in their sense of loneliness. The right companion at the right time can keep them from quitting midway.)}
\end{truyen}

\begin{baihoc}
Đi cùng đôi khi mạnh hơn thúc ép.
\textit{(Walking beside someone is sometimes stronger than pushing them.)}
\end{baihoc}

\begin{ghinhoanh}
\textbf{Short English to remember:}

Some students need company, not pressure.

\end{ghinhoanh}

\begin{tuhocnhanh}
1. Khi nản, em cần một lời hay một người ở cạnh hơn? \textit{(When discouraged, do you need words more, or someone's presence?)}

2. Ai là người khiến em thấy mình không một mình? \textit{(Who makes you feel less alone?)}
\end{tuhocnhanh}
\ngancach


\section{Thầy không bỏ em chỉ vì em đang xuống tinh thần.}
\begin{truyen}{Thầy không bỏ em chỉ vì em đang xuống tinh thần.}{Stay when students are low.}
\chuhoa{H}{ọc sinh nhớ rất lâu việc ai ở lại với mình lúc mình tệ nhất. Một người thầy không bỏ đi khi các em xuống tinh thần chính là một hình thức giáo dục sâu nhất.}
\textit{(Students remember for a long time who stayed with them at their worst. A teacher who does not walk away when their spirit drops is teaching at a very deep level.)}
\end{truyen}

\begin{baihoc}
Ở lại với học sinh lúc các em xuống tinh thần là một dạng dạy học bằng nhân cách.
\textit{(Staying with students when they are low is a form of teaching through character.)}
\end{baihoc}

\begin{ghinhoanh}
\textbf{Short English to remember:}

Stay when they are low.

\end{ghinhoanh}

\begin{tuhocnhanh}
1. Em có từng nhớ ai vì họ ở lại lúc em xuống tinh thần không? \textit{(Do you remember someone because they stayed when you were down?)}

2. Em muốn trở thành người như thế cho ai? \textit{(Who do you want to become that kind of person for?)}
\end{tuhocnhanh}
\ngancach


\section{Đừng chỉ hỏi em học đến đâu — hãy hỏi em còn muốn đi tiếp không.}
\begin{truyen}{Đừng chỉ hỏi em học đến đâu — hãy hỏi em còn muốn đi tiế...}{Check the heart, not only progress.}
\chuhoa{K}{hông phải lúc nào tiến độ cũng phản ánh sự thật bên trong. Có em vẫn làm bài đều nhưng trong lòng đã rất gần bờ bỏ cuộc.}
\textit{(Progress does not always reveal the inner truth. Some students keep doing the work even when they are very close to quitting inside.)}
\end{truyen}

\begin{baihoc}
Muốn giúp học sinh lâu dài, phải hỏi cả tiến độ lẫn tinh thần.
\textit{(To help students for the long term, ask about both progress and spirit.)}
\end{baihoc}

\begin{ghinhoanh}
\textbf{Short English to remember:}

Check the heart, not only the work.

\end{ghinhoanh}

\begin{tuhocnhanh}
1. Có khi nào em vẫn làm bài nhưng trong lòng muốn bỏ? \textit{(Have you ever kept working while secretly wanting to quit?)}

2. Em ước người lớn hỏi mình điều gì hơn? \textit{(What do you wish adults would ask you more often?)}
\end{tuhocnhanh}
\ngancach


\section{Một năm học không đáng để em mất luôn lòng tự trọng.}
\begin{truyen}{Một năm học không đáng để em mất luôn lòng tự trọng.}{Protect dignity.}
\chuhoa{K}{hi việc học được đặt cao hơn cả lòng tự trọng, học sinh sẽ học trong sợ hãi và xấu hổ. Điều đó có thể tạo ra kết quả ngắn hạn, nhưng để lại vết hằn lâu dài.}
\textit{(When study is valued more than dignity, students learn in fear and shame. That may create short-term results, but it leaves long-term wounds.)}
\end{truyen}

\begin{baihoc}
Giáo dục tốt không đạp lên lòng tự trọng để đổi lấy thành tích.
\textit{(Good education does not trample dignity for achievement.)}
\end{baihoc}

\begin{ghinhoanh}
\textbf{Short English to remember:}

Dignity matters in learning too.

\end{ghinhoanh}

\begin{tuhocnhanh}
1. Khi nào em thấy mình bị chạm vào lòng tự trọng vì chuyện học? \textit{(When have you felt your dignity hurt because of study?)}

2. Điều gì giúp em giữ được sự tự trọng của mình? \textit{(What helps you preserve your dignity?)}
\end{tuhocnhanh}
\ngancach


\section{Thầy mong em đi xa, nhưng thầy cũng mong em đi lành.}
\begin{truyen}{Thầy mong em đi xa, nhưng thầy cũng mong em đi lành.}{Go far, go whole.}
\chuhoa{Đ}{i xa mà kiệt quệ, đi nhanh mà rách lòng, đó không phải điều một người thầy thật sự mong. Đi lành nghĩa là còn nguyên nhân cách, sức khỏe và niềm tin.}
\textit{(Going far while exhausted, going fast while wounded inside, is not what a real teacher truly wants. Going well means keeping character, health, and hope intact.)}
\end{truyen}

\begin{baihoc}
Đi xa là tốt; đi lành mới bền.
\textit{(Going far is good; going whole is sustainable.)}
\end{baihoc}

\begin{ghinhoanh}
\textbf{Short English to remember:}

Go far, but go whole.

\end{ghinhoanh}

\begin{tuhocnhanh}
1. Em đang cố đi xa hay đang cố đi lành? \textit{(Are you only trying to go far, or also trying to go well?)}

2. Điều gì giúp em bền hơn trên đường học? \textit{(What helps you last longer on the learning road?)}
\end{tuhocnhanh}
\ngancach


\section{Hôm nay em còn ngồi đây là đã đáng quý rồi.}
\begin{truyen}{Hôm nay em còn ngồi đây là đã đáng quý rồi.}{Showing up matters.}
\chuhoa{C}{ó ngày học sinh không làm được gì nổi bật. Nhưng việc các em vẫn đến lớp, vẫn ngồi lại, vẫn chưa buông hoàn toàn, đã là một giá trị cần được nhìn nhận.}
\textit{(Some days students do nothing remarkable. But the fact that they still show up, still sit there, still have not fully given up, is already worth noticing.)}
\end{truyen}

\begin{baihoc}
Có mặt trong ngày khó cũng là một dạng can đảm.
\textit{(Showing up on a hard day is a form of courage.)}
\end{baihoc}

\begin{ghinhoanh}
\textbf{Short English to remember:}

Showing up is already brave.

\end{ghinhoanh}

\begin{tuhocnhanh}
1. Em có coi nhẹ việc mình vẫn cố đến lớp không? \textit{(Do you undervalue the fact that you still show up?)}

2. Em có thể tự ghi nhận mình điều gì hôm nay? \textit{(What can you acknowledge in yourself today?)}
\end{tuhocnhanh}
\ngancach


\section{Người thầy tốt không chỉ truyền bài — còn giữ học trò khỏi rơi khỏi đường học.}
\begin{truyen}{Người thầy tốt không chỉ truyền bài — còn giữ học trò khỏi ...}{Teaching is also keeping them in.}
\chuhoa{D}{ạy học không chỉ là chuyển kiến thức từ bảng vào vở. Có lúc, dạy học là giữ một học sinh ở lại với con đường của mình thêm một ngày, thêm một tuần, thêm một năm.}
\textit{(Teaching is not only moving knowledge from the board to the notebook. Sometimes teaching means helping one student stay on their path for one more day, one more week, one more year.)}
\end{truyen}

\begin{baihoc}
Giữ được học trò ở lại với việc học cũng là thành tựu của người thầy.
\textit{(Keeping students connected to learning is also an achievement of a teacher.)}
\end{baihoc}

\begin{ghinhoanh}
\textbf{Short English to remember:}

Teaching also means keeping students from falling away.

\end{ghinhoanh}

\begin{tuhocnhanh}
1. Theo em, người thầy tốt nhất đã giữ học sinh bằng điều gì? \textit{(In your view, how does the best teacher keep students going?)}

2. Em muốn nhớ câu nào nhất trong quyển này để đi tiếp? \textit{(Which sentence from this volume do you most want to remember to keep going?)}
\end{tuhocnhanh}
\ngancach


\section{Nếu hôm nay em chưa đứng dậy được hẳn, thầy vẫn chờ em ở ngày mai.}
\begin{truyen}{Nếu hôm nay em chưa đứng dậy được hẳn, thầy vẫn chờ em ở ...}{Tomorrow still waits.}
\chuhoa{C}{ó những ngày học sinh chỉ đủ sức để không ngã thêm. Thế cũng được. Người thầy tử tế hiểu rằng có những sự hồi phục cần thêm thời gian.}
\textit{(Some days students only have enough strength not to fall further. That is still okay. A kind teacher understands that some recoveries take more time.)}
\end{truyen}

\begin{baihoc}
Không phải ngày nào cũng là ngày bật dậy mạnh mẽ; có ngày chỉ cần chưa bỏ là đủ.
\textit{(Not every day is a dramatic comeback; some days it is enough simply not to give up.)}
\end{baihoc}

\begin{ghinhoanh}
\textbf{Short English to remember:}

Tomorrow is still waiting for you.

\end{ghinhoanh}

\begin{tuhocnhanh}
1. Nếu hôm nay em mới chỉ giữ được mình, điều đó có đáng quý không? \textit{(If today you only managed to hold yourself together, is that still valuable?)}

2. Em muốn đem điều gì của hôm nay sang ngày mai? \textit{(What do you want to carry from today into tomorrow?)}
\end{tuhocnhanh}
\ngancach


\section{Chỉ cần em còn muốn thử lại, thầy vẫn còn lý do để tin em.}
\begin{truyen}{Chỉ cần em còn muốn thử lại, thầy vẫn còn lý do để tin em.}{Trying again is enough reason for hope.}
\chuhoa{M}{ột học sinh chưa hoàn hảo nhưng còn muốn thử lại luôn đáng để đặt niềm tin. Người thầy tốt không đợi các em giỏi hẳn mới tin; thầy nhìn vào ý muốn đứng dậy của các em.}
\textit{(A student who is not perfect but still wants to try again is always worthy of trust. A good teacher does not wait until students are already strong to believe in them; he sees their desire to rise again.)}
\end{truyen}

\begin{baihoc}
Còn muốn thử lại là còn mầm hy vọng cần được gìn giữ.
\textit{(The desire to try again is a seed of hope worth protecting.)}
\end{baihoc}

\begin{ghinhoanh}
\textbf{Short English to remember:}

Trying again gives hope a reason.

\end{ghinhoanh}

\begin{tuhocnhanh}
1. Lần thử lại gần nhất của em là khi nào? \textit{(When was the last time you tried again?)}

2. Em muốn ai nhìn thấy sự cố gắng quay lại của mình? \textit{(Who do you want to notice your effort to return?)}
\end{tuhocnhanh}
\ngancach